%-----------------------------------------
% Chapter 4 - Social Impact and Sustainability
%-----------------------------------------

\chapter{SOCIAL IMPACT AND SUSTAINABILITY}

\section{SOCIAL IMPACT}

The proposed system has a great social impact as it shields individual users from social engineering attacks that target human trust, urgency, and authority instead of technical vulnerabilities. Such attacks include phishing, baiting, pretexting, whaling, and malicious URL attacks, which often lead to financial loss, identity theft, and unauthorized access to accounts. Many of these attacks go unnoticed until a great deal of damage has been done. The system will act as a smart first line of defense that runs continuously in the background, making it unnecessary for users to have advanced knowledge of cybersecurity. This is particularly important for senior citizens, students, and first-time internet users who are more prone to deceptive digital communication.

In addition to the personal security feature, the system enhances the overall cybersecurity of organizations and fosters digital trust. Social engineering attacks against executives and employees can result in fraudulent transactions, breaches of data, and reputation damage. With strong detection performance on all five attack classes, including strong recall on Whaling and strong F1-score on Pretexting in the incremental learning settings, the system assists in early threat detection, lowering the cost of response, and enhancing regulatory compliance. Furthermore, by leveraging open-source technologies and running efficiently on commodity hardware, the system assists in the democratization of cybersecurity, allowing small businesses, institutions of learning, and NGOs to benefit from cutting-edge cybersecurity protection.

\section{SUSTAINABILITY}

The proposed system is computationally efficient and sustainable. The Extreme Learning Machine calculates the output weights analytically without the need for iterative gradient descent computation, thus eliminating the need for GPU-intensive training. The Stacked Autoencoder is trained only once on the initial data and is left frozen for the rest of the incremental updates, thus eliminating the need for repeated deep learning training cycles. Furthermore, the Fuzzy-PGFDG feature calculations are performed through efficient matrix calculations that can be performed on standard CPU hardware. This is a significant move towards energy efficiency compared to the transformer model or the large language model approach.

Sustainability is further improved by the use of incremental adaptability and memory-efficient design. Unlike conventional machine learning systems, which have to be completely retrained when new types of attacks are discovered, the proposed incremental learning approach incorporates new knowledge without forgetting previously gained knowledge. This increases the useful life of the system. The memory cleanup protocol ensures that old feature matrices and models are deleted after each incremental learning update, regardless of the age of the system. This ensures that the system is scalable without increasing hardware costs, thus improving economic and environmental sustainability.